\section{Conclusion}
\label{sec:conclusion}

The system works and the design decision at the centre of it holds up. Keeping
the microcontroller authoritative made the client trivial to write and
impossible to cheat with, and the tests demonstrate the property directly rather
than arguing for it: every rejection case reads the board, issues a request that
must fail, and asserts the board is unchanged.

The decision with the largest practical payoff was separating the rules from the
hardware. Because \texttt{server/core/} has no HAL in it, the same sources run
on a host simulator, which meant the protocol and the game logic were finished
and tested before the board was involved at all. On a project with a fixed
deadline that converts bench time into laptop time, which is the cheaper of the
two.

\textbf{Limitations.} The server holds one game and one session, so it cannot
serve two clients at once; the state machine would need a session table rather
than a single \texttt{game\_t} to do so. Frames are capped at 64 bytes, which is
comfortable for this protocol but would need revisiting for a larger board. The
optional checksum detects corruption on the wire but is an XOR, so it is not
forgery-resistant, and integrity against a hostile client rests on server
authority instead. Transmission from the board is blocking with a timeout, which
is acceptable at 115\,200 baud and one short line per request but would not
scale to a chattier protocol.

\textbf{Further work.} The most natural extension is the AI opponent listed
among the bonus features, which fits the architecture without changing the
protocol: the server already owns the state, so it can generate a reply move
after a client move, and the existing suite could assert that it never plays
illegally. Beyond that, a session table would allow multiple concurrent games,
and a larger board would mostly be a matter of relaxing the frame limit and the
index parser.
